\chapter{PCB Hardware Design}
\label{chp:hardware}

\section{The Original Prototype}

The bearing and gradient algorithms were first realised on an earlier prototype built around a Teensy~4.1 microcontroller, which read the photodiodes correctly but was far too large, heavy, and power-hungry to fly on a micro-drone, and carried far more compute than the algorithms require.

The original board used larger, ready-made photodiode and operational-amplifier modules and provided space for sixteen photodiodes, more than the eight that Harry's prior research showed to be sufficient (Figure~\ref{fig:old-prototype}).

\begin{figure}[t]
    \centering
    \missingfigure[figwidth=0.6\textwidth]{Photograph of the original Teensy 4.1-based prototype board}
    \caption{The original Teensy 4.1-based prototype board.}
    \label{fig:old-prototype}
\end{figure}

This chapter presents a redesigned board that keeps the same sensing principle while meeting the size, weight, and power limits of the CrazyFlie platform, and compares the new design against this original throughout.

\section{Design Requirements}

The new board was designed to correct the shortcomings of the original prototype while meeting the CrazyFlie~2.1's hard limits: a total mass below 13 grams and a maximum diagonal footprint of 44\,mm.

It therefore uses eight photodiodes rather than sixteen, compact surface-mount photodiode and amplifier components in place of the larger ready-made modules, and offloads all computation to the CrazyFlie so that no power-hungry on-board microcontroller is required.

\section{Circuit Design}

Eight photodiodes are arranged in a ring and each connected to a transimpedance amplifier that converts the photodiode's current output to a measurable voltage signal.

The voltage outputs are digitised by an ADC module that communicates with the CrazyFlie's onboard processor over SPI, providing all eight channel readings at each sample.

The board includes optional first-order filter pads at each channel, allowing a low-pass or high-pass filter to be populated depending on the noise characteristics of the deployment environment.

\section{PCB Layout}

The eight photodiodes are uniformly distributed at 45\textdegree{} intervals around the full 360\textdegree{} circumference of the board to provide symmetric angular coverage.

A ground plane is included in the PCB layout to reduce the impact of electromagnetic interference generated by the drone's motors on the photodiode signal readings.

Before integration with the drone, the assembled board was verified to produce correct voltage readings on all eight channels using a standalone Teensy microcontroller.
